%%%%%%%%%%%%%%%%%%%%%%%%%%%%%%%%%%%%%%%%%%%%%%%%%%%%%%%%%%%
% EPFL report package, main thesis file
% Goal: provide formatting for theses and project reports
% Author: Mathias Payer <mathias.payer@epfl.ch>
%
% To avoid any implication, this template is released into the
% public domain / CC0, whatever is most convenient for the author
% using this template.
%
%%%%%%%%%%%%%%%%%%%%%%%%%%%%%%%%%%%%%%%%%%%%%%%%%%%%%%%%%%%
\documentclass[a4paper,11pt,oneside]{report}
% Options: MScThesis, BScThesis, MScProject, BScProject
\usepackage[MScThesis,lablogo]{EPFLreport}
\usepackage{xspace}

% added packages
\usepackage{amsmath}
\usepackage{cleveref}
\usepackage{float}
\usepackage{tabularx}
\usepackage{booktabs}
\usepackage{xltabular}
\usepackage[version=4]{mhchem}
\usepackage[colorinlistoftodos]{todonotes} % for the citing todo
\newcommand{\needcite}[1]{\todo[color=red!40]{CITE: #1}}

\title{Tritium extraction and transport from molten salt breeding experiments}
\author{Hippolyte Monroe}
\supervisor{Dr. Rémi Delaporte-Mathurin}
\adviser{Dr. Mathieu Hursin}
%\coadviser{Second Adviser}
\expert{The External Reviewer}

% custom commands
\newcommand{\sysname}{FooSystem\xspace}
\newcommand{\PT}{\ensuremath{P_{T_2}}}
\newcommand{\CT}{\ensuremath{c_{T_2}}}
\newcommand{\JT}{\ensuremath{J_{T_2}}}
\newcommand{\YT}{\ensuremath{y_{T_2}}}

\begin{document}
\maketitle
\listoftodos
\makededication

\acknowledgments{
my ack 
\par
I would also like to acknowledge the support of the Zeno Karl Schindler foundation.

}
\makeacks

\begin{abstract}
The \sysname tool enables lateral decomposition of a multi-dimensional
flux compensator along the timing and space axes.

The abstract serves as an executive summary of your project.
Your abstract should cover at least the following topics, 1-2 sentences for
each: what area you are in, the problem you focus on, why existing work is
insufficient, what the high-level intuition of your work is, maybe a neat
design or implementation decision, and key results of your evaluation.
\end{abstract}

\begin{frenchabstract}
For a doctoral thesis, you have to provide a French translation of the
English abstract. For other projects this is optional.
\end{frenchabstract}

\maketoc

\chapter*{Notations}
\addcontentsline{toc}{chapter}{Notations}
\markboth{Notations}{Notations}

\renewcommand{\arraystretch}{1.4}
\setlength{\tabcolsep}{6pt}

\begin{xltabular}{\linewidth}{@{}l X l@{}}
\multicolumn{3}{@{}l}{\textit{\textbf{Latin symbols}}} \\
\midrule
\endfirsthead

\multicolumn{3}{@{}l}{\textit{\textbf{(continued)}}} \\
\midrule
\endhead

\midrule
\multicolumn{3}{r@{}}{\textit{continued on next page}} \\
\endfoot

\endlastfoot

$A$          & Tank (column) base area                     & m$^2$ \\
$a$            & Specific interfacial area                  & m$^{-1}$ \\
$c_{T_2}$      & Tritium concentration in liquid            & mol\,m$^{-3}$ \\
$d_b$          & Sauter mean bubble diameter                       & m \\
$d_n$          & Sparger nozzle diameter                    & m \\
$D_l$          & Tritium diffusivity in liquid              & m$^2$\,s$^{-1}$ \\
$E$            & Phase axial dispersion coeff.              & m$^2$\,s$^{-1}$ \\
$h_l$          & Liquid-side mass transfer coeff.           & m\,s$^{-1}$ \\
$H$            & Liquid column height                       & m \\
$K_s$          & Solubility constant                        & mol\,m$^{-3}$\,Pa$^{-1}$ \\
$M_{\mathrm{He}}$ & Molar mass of helium                    & kg\,mol$^{-1}$ \\
$\dot n_g$     & Molar gas flow rate                        & mol\,s$^{-1}$ \\
$N_\alpha$          & Total tritium inventory in phase $\alpha$                 & mol \\
$N_n$          & Number of sparger nozzles                  & -- \\
$P$            & Pressure                                   & Pa \\
$P_{\mathrm{top}}$ & Pressure at column top                & Pa \\
$P_{\mathrm{bot}}$ & Hydrostatic pressure at column bottom & Pa \\
$P_{T_2}$      & Tritium partial pressure in gas            & Pa \\
$u_g$          & Superficial gas velocity                   & m\,s$^{-1}$ \\
$v_g$       & Bubble true (interstitial) velocity                   & m\,s$^{-1}$ \\
$\dot V_g$     & Volumetric gas flow rate                   & m$^3$\,s$^{-1}$ \\
$z$            & Axial coordinate                           & m \\

\addlinespace[1em]
\multicolumn{3}{@{}l}{\textit{\textbf{Greek symbols}}} \\
\midrule
$\varepsilon$ & Volumetric phase fraction                   & -- \\
$\mu_l$       & Liquid dynamic viscosity                    & Pa\,s \\
$\nu_l$       & Liquid kinematic viscosity                  & m$^2$\,s$^{-1}$ \\
$\rho_g$      & Gas-phase density                           & kg\,m$^{-3}$ \\
$\rho_l$      & Liquid-phase density                        & kg\,m$^{-3}$ \\
$\Delta\rho$  & Liquid--gas density difference              & kg\,m$^{-3}$ \\
$\sigma_l$    & Liquid surface tension                      & N\,m$^{-1}$ \\

\addlinespace[1em]
\multicolumn{3}{@{}l}{\textit{\textbf{Subscripts and abbreviations}}} \\
\midrule
$0$    & Tank bottom (gas inlet)             & \\
$b$    & Bubble                              & \\
$g$    & Gas phase                           & \\
$l$    & Liquid phase                        & \\
$i$    & Liquid-gas interface                & \\
$n$    & Nozzle                              & \\
$t$    & Tank                                & \\
$T_2$  & Molecular tritium                   & \\
ARD    & Advection--reaction--dispersion     & \\
FLiBe  & Molten salt Li$_2$BeF$_4$           & \\
$\mathrm{TBR}$ & Tritium breeding ratio       & -- \\

\end{xltabular}

%%%%%%%%%%%%%%%%%%%%%%
\chapter{Introduction}
%%%%%%%%%%%%%%%%%%%%%%

% The introduction is a longer writeup that gently eases the reader into your
% thesis~\cite{dinesh20oakland}. Use the first paragraph to discuss the setting.
% In the second paragraph you can introduce the main challenge that you see.
% The third paragraph lists why related work is insufficient.
% The fourth and fifth paragraphs discuss your approach and why it is needed.
% The sixth paragraph will introduce your thesis statement. Think how you can
% distill the essence of your thesis into a single sentence.
% The seventh paragraph will highlight some of your results
% The eights paragraph discusses your core contribution.

% This section is usually 3-5 pages.

\section{The tritium isotope}
% what is tritium (main characteristics)
Tritium is a radioactive isotope of hydrogen, its nucleus is made of 1 proton and 2 neutrons and is therefore denoted \ce{^3H} or \ce{T}. It disintegrates into \ce{^3He} through beta decay with a half life of 12.3 years. 
% where do we find it
Because of its short half life and its absence of any decay chain, tritium can't be naturally found in concentrated quantities unlike other radioactive isotopes (e.g. uranium-238, radon-222). It is unintentionally produced: 
\begin{itemize}
    \item naturally, from highly energetic cosmic neutrons interacting with atmospheric nitrogen (\ce{^14_7N(n,T)^12_6C} reaction, contribution estimated at 0.2 kg/year)
    \item as a product of the fission reactions in all fission reactors, with a low yield ($10^{-4}$ to $10^{-6}$ T/fission event, depending on the fissile nuclide and the neutron energy) \needcite{IAEA database} %https://www-nds.iaea.org/relnsd/vcharthtml/VChartHTML.html
\end{itemize}

% do we want tritium ?
Even though beta radiations don't penetrate the skin, tritium can substitute to hydrogen in water to form \ce{HTO}, which if spread in the environment can contaminate the food chain and become a radiological hazard. For that reason, it is an undesirable byproduct in the fission industry. % -> no in fission

Still, tritium finds a use as a fuel in nuclear fusion: because of its favorable cross section for fusion reactions with deuterium (\ce{^2H} or \ce{D}), it was identified as the most efficient fuel for fusion devices \needcite{ITER physical basis report ?}. Historically, nuclear fusion and therefore tritium have been serving military purposes, in nuclear weapons. But just like nuclear fission found a peaceful use in nuclear reactors for energy production, controlled nuclear fusion offers a promising path to carbon-free energy. Compared to its fission counterpart, a fusion power plant (FPP) would present the advantage of lower and shorter radiotoxicity of its byproducts, increased safety and abundance of fuel. Controlled fusion has therefore been the subject of intense research for decades, and is now entering a fast growing phase: in 2026, the Fusion Industry Association counted 56 private fusion companies for a total funding of \$14.2bn worldwide, which is \$4.5bn more than in 2025\cite{noauthor_2026_2026}. Among these companies, 40 of them envision deuterium-tritium (DT) as fusion fuel. Deuterium is present in small proportions in seawater (\needcite{how many ppm ?}) and its extraction is a well known process as evidenced by the fleet of 44 Pressurized Heavy Water Reactors (PHWR) operating worldwide. Tritium fuel supply on the other hand is a much bigger concern due to itsp scarcity, stressing the need for industrial tritium production. % -> yes in fusion

\todo{insert representation of H and T nucleus + isotope chart along with energy barrier}
% elaborate on why do we want fusion power plants ? Dedicate a subsection ?

\section{Tritium production}
\subsection{Main pathways}
% how can we create tritium 
A typical FPP would require important amounts of tritium, estimated at $\sim$55 kg/year per GW of fusion power \cite{abdou_physics_2020}. The current world inventory of $\sim$20 kg of tritium is largely insufficient and halves every 12 years due to radioactive decay if not renewed. It is then necessary to intentionally produce tritium via nuclear transmutation:
\begin{align}
%include reaction with B10 ? -> in PWR 
    \ce{^2H + n &-> T + \gamma} \label{eq:D2_capture}\\
    \ce{^6Li + n &-> T + \alpha} + 4.8~\mathrm{MeV} \label{eq:Li6_capture}\\
    \ce{^7Li + n &-> T + \alpha} + \mathrm{n} - 2.5~\mathrm{MeV} \label{eq:Li7_capture}
\end{align}

As of now, the existing commercially available tritium is produced as a byproduct of heavy water moderated reactors through neutron capture (\cref{eq:D2_capture}), specifically in CANada Deuterium Uranium,(CANDU) reactors (estimated at 130 g of tritium/reactor/year) \needcite{https://www.neimagazine.com/analysis/answering-the-big-tritium-question/?cf-view}. But the small cross section of this reaction (0.5 millibarns for thermal neutrons) makes it unsuitable for large scale production. \Cref{eq:Li6_capture} on the other hand  has a 1000~barns thermal neutron cross section \needcite{JANIS database (https://www.oecd-nea.org/janisweb/book/neutrons/H2/MT102/renderer/3458)}, it is used for military tritium production and is also the privileged path to tritium breeding for controlled fusion. Industrial tritium production would therefore rely on thermal neutron capture in a lithium-6 rich material, the cross section for \cref{eq:Li7_capture} being marginal (300 millibarns, for fast neutrons). 

% where do we get the neutrons from
Aforementioned transmutation reactions all need neutrons as a reactant. Neutrons can be found in fission reactors, but producing tritium there poses the problem of introducing large amounts of tritium in an environment where it is usually unwanted because of its handling difficulties and the associated radiological concern, and is anyway constrained to relatively low tritium amounts to leave enough neutrons for fission to self-sustain ($\sim$1 kg/reactor/year using Tritium Producing Burnable Absorbers Rods)\needcite{NRC?}. For those reasons, tritium production in fission reactors is restricted to military programs.

DT fusion itself produces fast neutrons (\cref{eq:DT_fusion}), it naturally follows that a FPP can act as a neutron source that \textit{breeds} its own fuel. Contrary to external sources such as fission reactors or neutron beams, this approach offers the advantage that the available neutron flux scales with the amount of fuel needed and thus can sustain operation of an arbitrarily large fleet of FPP. This makes breeding the most scalable and economical path to industrial tritium fuel production. 

\begin{align}
    \ce{D + T -> \alpha}\ \mathrm{(3.5 MeV) + n \ (14 MeV)} \label{eq:DT_fusion}
\end{align}
    

\todo{insert plot of neutron cross sections of those reactions} % use Remi's neutron XS plotter: https://github.com/openmc-data-storage/isotope-xs-plotter

\subsection{The breeding blanket}
% how is breeding engineered ?
In practice, a FPP would produce tritium in a \textit{breeding blanket} that surrounds the fusion core, be it inertially or magnetically confined. For \cref{eq:Li6_capture,eq:Li7_capture} to occur, the breeding blanket is made of a lithium-containing material, chosen among:
\begin{itemize}
    \item solid ceramics: lithium orthosilicate (\ce{Li4SiO4}), lithium metatitanate (\ce{Li2TiO3})
    \item liquid metals: Lithium-Lead eutectic (\ce{Pb_{83}Li_{17}}), melting at 235°C
    \item molten salts: FLiBe (\ce{2LiF-BeF2}), FLiNaK (\ce{LiF-NaF-KF}), melting at $\sim$450°C
\end{itemize}

% material choice: tritium production
From the perspective of tritium production optimization, the choice of breeding material is driven by: 
\begin{itemize}
    \item a high macroscopic neutron capture cross section: higher Li density, \ce{ ^6Li} enrichment is favorable
    \item presence of a neutron moderator: the \ce{^6Li} reaction occurs with thermal neutrons, therefore light elements are favorable for their moderating capability
    \item presence of a neutron multiplier: Pb and Be containing materials readily incorporate a neutron multiplier ((n,2n) reaction)
\end{itemize}

% the other roles of the BB
Beyond the tritium production aspect, a breeding blanket must shield the 14 MeV neutrons produced by the fusion core, i.e. 80\% of the fusion energy (see \cref{eq:DT_fusion}) dissipates in it as heat. In that regard, liquid materials offer the advantage of being both tritium breeding materials and coolant. Molten salts in particular have a heat capacity superior to that of LiPb and don't conduct electricity which makes them favorable for compact and magnetically confined FPP designs, such as the Affordable-Robust-Compact (ARC) tokamak concept. 


\section{Challenges of the tritium fuel cycle}
The fusion fuel cycle is the set of processes that supplies tritium to the plasma, divided in two independent closed loops: % "independent": outer loop doesn't have the output of inner loop as input, contrary to what the "inner" and "outer" denominations suggest
\begin{itemize}
    \item \textit{inner fuel cycle}: recovery of what the plasma does not burn. It is characterized by the fueling efficiency (performance of the fuel injection system, i.e. how much fuel reaches the fusion core) and the tritium burn fraction TBF (performance of the fusion core, i.e. how much tritium fuses among the one that reached the core). The product of these two quantities isn't expected to exceed 5\% \cite{abdou_physics_2020}
    \item \textit{outer fuel cycle}: newly bred tritium to replace what was consumed. It is characterized by the Tritium Breeding Ratio (TBR), defined as the ratio of tritium bred per tritium consumed. 
\end{itemize}

\todo{insert schematic of fuel cycle, from Meschini 2023 ARC fuel cycle paper ?}

\subsection{Self-sufficiency}
% definition of the scientific and engineering TBR, tritium self-sufficiency of a FPP
% include doubling time considerations, achievable and required TBR ? (Ferry 2023, Abdou 2020)
For the fusion power industry to be able to scale, individual FPPs have to produces at least as much tritium as they consumes, i.e. optimize their outer fuel cycle to achieve a sufficiently high TBR. For explanatory purposes. a distinction can be made between the "scientific" TBR (\cref{eq:TBR_sci}), which only measures tritium being born in the breeding material, and the "engineering" TBR (\cref{eq:TBR_eng}) which measures the tritium that is actually retrieved and can be injected in the fusion core. 
\begin{align}
    \mathrm{TBR_{sci}}&= \frac{\mathrm{T_{bred}}}{\mathrm{T_{burnt}}} \label{eq:TBR_sci} \\
    \mathrm{TBR_{eng}} &= \frac{\mathrm{T_{retrieved}}}{\mathrm{T_{burnt}}} =  \mathrm{TBR_{sci}} \cdot \eta  \label{eq:TBR_eng}\\
    \eta &= \frac{\mathrm{T_{retrieved}}}{\mathrm{T_{bred}}}<1 \label{eq:eta}
\end{align}

Where $\eta$ accounts for losses through the outer fuel cycle. 

$\mathrm{TBR_{sci}}$>1 is necessary, but it is not sufficient. A self sufficient FPP needs to achieve $\mathrm{TBR_{eng}}$>1 by maximizing both its components:
\begin{itemize}
    \item $\mathrm{TBR_{sci}}$ is purely neutronics dependent, it is optimized by use of \ce{^6Li} enrichment and by maximizing its thickness as well as the solid angle it covers around the fusion core;
    \item $\eta$ is optimized by reducing permeation losses throughout the fuel cycle (e.g. by means of dedicated structural materials and coatings) and optimizing tritium extraction from the breeder.
\end{itemize}  

\subsection{Tritium handling}
% tritium is very mobile
As isotopes of the same element share same electron structure, tritium chemically behaves like hydrogen (notwithstanding a few nuances in their kinetic properties arising from their mass difference). Tritium forms \ce{T2} gas and can substitutes with hydrogen to form tritiated water \ce{HTO}. It can dissolve as atomic \ce{T} in metallic materials, as a consequence of hydrogen being a small but non-inert element, enabling a 3-step process called \textit{permeation}: 
\begin{enumerate}
    \item tritium bonds with the conduction band electrons at the metal surface (dissociation of \ce{T2} into adsorbed \ce{T})
    \item it then fits into interstitial sites of the lattice (high \textit{solubility} in metals) and hop between them over a low energy barrier (high \textit{diffusivity})
    \item  until it reaches the opposite free surface where it becomes \ce{T2} gas again (recombination and desorption)
\end{enumerate}
Permeation is characterized by \textit{permeability} = diffusivity $\times$ solubility, and makes tritium especially hard to contain no matter how leak-tight and defect-free the container is. 

Permeation represents a significant part of outer fuel cycle losses (i.e. degrades $\eta$) and can pose a radiological issue. This is especially the case in liquid breeders that need to operate above a certain melting point, since permeability increases with temperature. 

\subsection{Tritium extraction}
On top of directly retrieving tritium, maximizing the efficiency of the tritium extraction system (TES) has the indirect positive effect of reducing the amount of tritium available for permeation, therefore decreasing tritium permeation losses in the outer fuel cycle. A performant TES is therefore needed to achieve a sufficiently high $\eta$. 

Several TES technologies exist, with various maturities:
\begin{itemize}
    \item Vacuum permeators:
    \item Gas-liquid contactors: 
    \item Vacuum sieve tray: 
\end{itemize}

\section{State of the art and existing gaps}

Both inner and outer fuel cycle are the subject of intense research, with the aim not to lose tritium and maximize its production. 

On the outer fuel cycle side, the main focus is to better understand tritium handling and extraction, in other words to be able to predict and maximize $\mathrm{TBR_{eng}}$. In the context of the International Thermonuclear Experimental Reactor (ITER), the OFC aspect of ceramic and liquid metal based breeders was investigated in several studies\needcite{ITER references}. 

There is however large uncertainties when it comes to molten salt based breeders, due to the challenges involved in experimenting with them: high temperatures, corrosion, Be toxicity and general tritium handling. This gap is addressed by the Liquid Immersion Blanket:Robust Accountancy (LIBRA) experiment, which main goal is to demonstrate accurate prediction of tritium breeding (the $\mathrm{TBR_{sci}}$ component) and extraction (the $\eta$ component), in a FLiBe salt breeder and DT fusion conditions (14 MeV neutrons)\cite{ferry_libra_2023}. 

Early results from this program are encouraging. The BABY 1L campaign, conducted on litre-scale salt volumes, showed good tritium recovery without any dedicated extraction system: tritium was collected simply by letting it diffuse through the salt to the free surface, or permeate through the container walls \cite{delaporte-mathurin_baby_2026}. At this scale, passive transport is sufficient.

This approach does not survive scaling, as the characteristic time for tritium to diffuse out of the salt grows quadratically with the salt dimensions. The forthcoming LIBRA Pi experiment will operate on a salt volume of $\sim$100\,L, its characteristic diffusion time is $\frac{L^2}{D}=\frac{(30\ \mathrm{cm})^2}{5\cdot10^{-9} \ \mathrm{m^2/s}}\approx200 \ \mathrm{days}$ if no convective transport is considered. Such a timescale is incompatible with the experimental iteration rate the program requires, and an active extraction system therefore becomes necessary.

Gas sparging was identified as a candidate to fill this role. The principle is to bubble an inert gas (typically helium) through the liquid breeder: tritium transfers from the liquid to the rising bubbles across the gas-liquid interface and is carried out of the tank, while the motion of the bubbles simultaneously promotes advective tritium transport within the liquid itself. The process is affordable, mechanically simple, and well documented in the chemical engineering literature \needcite{BCR / sparging chemical engineering references}. 

\section{Problem definition}
% roadmap: 1 sentence/chapter 

The viability of gas sparging as a tritium extraction system for LIBRA Pi has not been established yet. It is the object of the present modelling work, focusing on three main questions:
\begin{itemize}
    \item \textit{What extraction timescales could a sparging system achieve in LIBRA Pi ?}
    \item \textit{How can the extraction speed be optimized ?}
    \item \textit{What are the limitations of this prediction and how can it be improved ?}
\end{itemize}

\Cref{sec:background} will introduce the prior work on gas sparging upon which this work is built. 

\Cref{sec:sparging_model} will present the general gas sparging species extraction model developped in this work, beginning with the description of the underlying mathematical model (\cref{sec:sparging_maths}), following with its practical implementation (\cref{sec:sparging_impl}), then presenting the engineering insights that can be learnt from this model (\cref{sec:sparging_expl}). 

\Cref{sec:LIBRAPi_app} will apply the developed model to the LIBRA Pi experiment, first focusing on extraction time prediction (\cref{sec:LIBRAPi_extraction}), which will then be used to inform a study of tritium permeation in LIBRA Pi (\cref{sec:LIBRAPi_permeation}). 

%%%%%%%%%%%%%%%%%%%%
\chapter{Background}
%%%%%%%%%%%%%%%%%%%%
\label{sec:background}
% The background section introduces the necessary background to understand your
% work. This is not necessarily related work but technologies and dependencies
% that must be resolved to understand your design and implementation.

% This section is usually 3-5 pages.

% the literature review goes here, not in the intro
Helium sparging was used in the Molten Salt Reactor Experiment (MSRE), primarily intended for Xe-135 poison removal but also impacting the tritium distribution\cite{briggs_calculation_1970}. 

%%%%%%%%%%%%%%%%
\chapter{Gas sparging model}
\label{sec:sparging_model}
%%%%%%%%%%%%%%%%

% Introduce and discuss the design decisions that you made during this project.
% Highlight why individual decisions are important and/or necessary. Discuss
% how the design fits together.

% This section is usually 5-10 pages.

\section{Mathematical framework}
\label{sec:sparging_maths}
\subsection{Mass transfer}
% expose diffusivity, boundary layer model, then advective transport 
\subsection{Fluid dynamics}
% expose correlations

\subsection{Advection-reaction-dispersion model}
\begin{itemize}
    \item Liquid: 
    \begin{equation}
    \label{eq:CT2}
        \varepsilon_l \frac{\partial c_{T_2}}{\partial t} =  \underbrace{\varepsilon_l E_{l} \frac{\partial ^2 c_{T_{2}}}{\partial z^2}}_\text{liquid dispersion} + \underbrace{S_{T_2}}_\text{tritium source} - \underbrace{a\,J_{T_2}}_{\text{mass transfer}}
    \end{equation}
    \item Gas: 
    \begin{equation}
    \label{eq:PT2}
        \varepsilon_g \frac{1}{RT} \frac{\partial P_{T_2}}{\partial t} = \underbrace{-\frac{1}{RT}\frac{\partial (\,u_g\,P_{T_2})}{\partial z}}_\text{gas advection} + \underbrace{\varepsilon_g E_{g} \frac{1}{RT}\frac{\partial ^2 P_{T_{2}}}{\partial z^2}}_\text{gas dispersion} + a\,J_{T_2}
    \end{equation}
\end{itemize}
\subsection{Boundary conditions}
\subsection{Assumptions of the model}

\section{Numerical implementation}
\label{sec:sparging_impl}
% The implementation covers some of the implementation details of your project.
% This is not intended to be a low level description of every line of code that
% you wrote but covers the implementation aspects of the projects.

% This section is usually 3-5 pages.

\subsection{Python package structure}
% just a quick figure of the packages dependencies, API
\subsection{Numerical scheme}
\subsection{Fluid dynamics parameters resolution}
% parameters dependance graph
% graph resolution algorithm 

\subsection{Verification}
\subsubsection{Method of manufactured solution}
\subsubsection{Numerical error estimate}
% convergence plot: quantify convergence without tau at that point ? (not already introduced)
% add a MMS verification


\section{Model exploration}
\label{sec:sparging_expl}
% all that was learn by playing around with the model
\subsection{Effect of mixing}
\subsubsection{Continuous stirred tank approximation}
\subsubsection{Plug flow reactor approximation}

\subsection{Limiting regimes}
% plot error vs Pi 
\subsubsection{Mass transfer limited}
% sobol at low Pi -> Ks negligible
\subsubsection{Partial pressure limited}
% sobol at higher Pi -> Ks not negligible
\subsection{Analytical 0D model}
\subsubsection{Derivation}
Let us consider a single He bubble $b$, immersed in a static liquid of tritium concentration $c_l$. The evolution of the tritium inventory $n_{T_2,b}$ in the bubble is governed by:
\begin{equation}
\label{eq:nb}
    \frac{d\,n_{T_2,b}}{dt} = A_b\,\JT
\end{equation}

Writing it in terms of tritium partial pressure by using the ideal gas law and expanding the right term, \cref{eq:nb} becomes:

\begin{equation}
\label{eq:Pb_ODE}
    \frac{V_b}{RT}\frac{d\,\PT}{dt} = A_b\,h_l\,(c_l-K_s\PT)
\end{equation}

We are assuming a spherical bubble of diameter $d_b$, the specific interfacial area of the bubble is $a_b=\frac{A_b}{V_b}=\frac{6}{d_b}$ (not to be confused with the specific interfacial area of the gas phase as a whole $a=\varepsilon_ga_b$).
Solving \cref{eq:Pb_ODE} yields:
\begin{equation}
\label{eq:Pb_full}
    \PT(t)=\frac{c_l}{K_s}\,(1-e^{-a_b\,R\,T\,h_l\,K_s\,t})
\end{equation}

The speed $v_b$ of the bubble, assumed constant, couples the time it spends in the system to its position: $t=\frac{z}{v_b}$. We can express the member of the exponential term in \cref{eq:Pb_full} as a dimensionless number $\Pi$ defined as:
\begin{equation}
\label{eq:Pi_def}
    \Pi(z)\equiv \frac{a_b\,R\,T\,h_l\,K_s}{v_g}\,z
\end{equation}

$\Pi$ can be interpreted as "how close the bubble is to reaching $\frac{c_l}{K_s}$, i.e. being at equilibrium with the surrounding liquid". 
\begin{itemize}
    \item When $1-e^{\Pi(z)}\approx0$: the bubble can be considered free of any tritium, the driving force for tritium exchange \JT is therefore maximal, limited only by the liquid side mass transfer $h_l$, i.e. the ability of tritium to diffuse down the concentration gradient, through the liquid film up to the gas phase. 
    \item When $1-e^{\Pi(z)}\approx1$: the bubble is saturated with tritium, the interfacial flux \JT drops as there is no concentration gradient through the liquid film to drive it, since the mass transfer process is throttled by the partial pressure in the bubble being at equilibrium with the liquid phase.
\end{itemize}



(i.e. it can leave the boundary of the system eventually)
\subsubsection{Extension to partial pressure limited regime}
% fit e(Pi) with polynomial to get tau(Pi) 

%%%%%%%%%%%%%%%%%%%%%%%%
\chapter{Application to the LIBRA Pi tritium breeding experiment}
%%%%%%%%%%%%%%%%%%%%%%%%
\label{sec:LIBRAPi_app}

\section{Prediction of the extraction time}
\label{sec:LIBRAPi_extraction}
\subsection{Dependence on operating parameters}
\subsubsection{Sensitivity analysis}
% sobol plot of operating parameters
\subsubsection{Operating point}
% tau vs (n,T) map
% add cost function to give optimal region ?
\subsection{Uncertainty on transport parameters}
% given Pi number (previous chapter), can Ks be neglected
% pessimistic vs optimistic scenario

\section{Tritium permeation study}
\label{sec:LIBRAPi_permeation}
\todo{study considers we are in diffusion-limited regime, but check W>>1 from Abdou 2020}
\subsection{Tritium transport model}
\subsubsection{Geometry}
\subsubsection{Numerical scheme}
\subsection{Grid refinement study}

% safety factor SF=1.25 as recommanded for GCI on 3 or more grids
% we use P1 linear lagrange elements which theoretically give an L2 norm of the error in O(h^2) for the field themselves and O(h) for their gradient. Bounday fluxes are obtained from gradient so the results p=1-1.3 are consistent with the theory 

\subsection{Results}
\subsubsection{Impact of sweep gas composition}
\subsubsection{Impact of recombination constant}
\subsubsection{Effect of tritium extraction kinetics}


%%%%%%%%%%%%%%%%%%%%
\chapter{Evaluation}
%%%%%%%%%%%%%%%%%%%%

% In the evaluation you convince the reader that your design works as intended.
% Describe the evaluation setup, the designed experiments, and how the
% experiments showcase the individual points you want to prove.

% This section is usually 5-10 pages.


%%%%%%%%%%%%%%%%%%%%%%
\chapter{Related Work}
%%%%%%%%%%%%%%%%%%%%%%

The related work section covers closely related work. Here you can highlight
the related work, how it solved the problem, and why it solved a different
problem. Do not play down the importance of related work, all of these
systems have been published and evaluated! Say what is different and how
you overcome some of the weaknesses of related work by discussing the 
trade-offs. Stay positive!

This section is usually 3-5 pages.


%%%%%%%%%%%%%%%%%%%%
\chapter{Conclusion}
%%%%%%%%%%%%%%%%%%%%

In the conclusion you repeat the main result and finalize the discussion of
your project. Mention the core results and why as well as how your system
advances the status quo.

\cleardoublepage
\phantomsection
\addcontentsline{toc}{chapter}{Bibliography}
\printbibliography

% Appendices are optional
% \appendix
% %%%%%%%%%%%%%%%%%%%%%%%%%%%%%%%%%%%%%%
% \chapter{How to make a transmogrifier}
% %%%%%%%%%%%%%%%%%%%%%%%%%%%%%%%%%%%%%%
%
% In case you ever need an (optional) appendix.
%
% You need the following items:
% \begin{itemize}
% \item A box
% \item Crayons
% \item A self-aware 5-year old
% \end{itemize}

\end{document}